% CLAUDE REPORT::WCDERIV-0001::2026-11-11
\documentclass[11pt]{article}
\usepackage{amsmath,amssymb,amsthm}
\newtheorem{theorem}{Theorem}\newtheorem{proposition}[theorem]{Proposition}
\theoremstyle{remark}\newtheorem{remark}[theorem]{Remark}
\title{Worst-case minimum derivation cost $\overline W_R(N)$:\\ status of the bridge $\mathrm{SAT}\in\mathrm P\Rightarrow \overline W_R(N)\le\mathrm{poly}(N)$}
\author{Claude routine WCDERIV-0001}\date{2026-11-11}
\begin{document}\maketitle
\paragraph{Status.} P versus NP remains \textbf{OPEN}. This report gives no lower bound for arbitrary
algorithms. It classifies the proposed bridge and finds where it fails.

\section{Object}
For a fixed sound and complete inference system $R$ (Cook--Reckhow: a polynomial-time checker),
\[
\overline W_R(N)=\max_{F\in\mathrm{UNSAT},\,|\langle F\rangle|\le N}\ \min_{\pi:F\vdash_R\bot}\operatorname{cost}_R(\pi).
\]
Taking the minimum over all derivations is the same as nondeterministic proof length. Strategy search is not
charged, which is legitimate for a lower bound. The multisign DAG-hypergraph reduction calculus is an instance:
its ``minimal reduced logical state'' for $F\in\mathrm{UNSAT}$ is always $\bot$, so all the content is in
the cost term. If the calculus can introduce new hyperedges or variables (extension), it is at least as strong
as Extended Frege (EF). Otherwise it is a resolution-type or bounded-width calculus.

\section{Results}
\begin{proposition}[PROVED/EXACT; Cook--Reckhow 1979]\label{p:cr}
If $\overline W_R(N)\le\mathrm{poly}(N)$ for some $R$, then $\mathrm{NP}=\mathrm{coNP}$. So a superpolynomial
$\overline W_R$ is necessary, but it is not sufficient for $\mathrm P\ne\mathrm{NP}$. The bridge is the
load-bearing arrow.
\end{proposition}

\begin{theorem}[PROVED/EXACT: reformulation for weak $R$]\label{t:reform}
Let $R$ be any system with an unconditional superpolynomial lower bound for $\overline W_R$, such as
tree-like or dag resolution (Haken 1985: $\mathrm{PHP}^{n+1}_n$; Urquhart 1987: Tseitin), cutting planes,
or bounded-depth Frege. Then
\[
\bigl[\mathrm{SAT}\in\mathrm P\Rightarrow\overline W_R(N)\le\mathrm{poly}(N)\bigr]\iff \mathrm P\ne\mathrm{NP}.
\]
\end{theorem}
\begin{proof}The consequent is false, so the implication holds exactly when its antecedent is false.\end{proof}
So for every system where the agents can actually certify hard instances, proving the bridge is exactly as hard
as P$\ne$NP. In world-ledger terms it is an equivalent, not progress. Any hard-instance work for such $R$
moves only the consequent, which is already settled.

\begin{theorem}[PROVED/EXACT: relativization firewall]\label{t:rel}
Let $R$ be as in Theorem~\ref{t:reform}. Let $A$ be a PSPACE-complete oracle, so $\mathrm{SAT}\in\mathrm P^A$.
The relativized bridge ``$\mathrm{SAT}\in\mathrm P^A\Rightarrow\overline W_R\le\mathrm{poly}$'' is
\emph{false}, because $\overline W_R$ is defined without oracles and its lower bound holds in every world.
Hence any proof of the bridge must use the oracle-free code of the hypothetical SAT algorithm in an
essential way. ``Simulate the algorithm's steps as black-box reduction moves'' relativizes and cannot work.
\end{theorem}

\begin{proposition}[KNOWN THEOREM; Cook 1975, Kraj\'{\i}\v{c}ek--Pudl\'ak]\label{p:pv}
If a polynomial-time algorithm $M$ decides SAT and the theory $\mathrm{PV}$ (equivalently $S^1_2$ for this
$\forall\Sigma^b_1$ statement) proves its soundness on UNSAT outputs, then EF has polynomial-size refutations
of all UNSAT CNFs. So $\overline W_{\mathrm{EF}}\le\mathrm{poly}$.
\end{proposition}
This is the strongest bridge in the literature, and it is \emph{weaker} than the requested one. The missing
case is a correct algorithm whose correctness PV does not prove. The bridge does not fail because of strategy
search; the free-optimal-sequence principle is sound. It fails because an algorithm's ``UNSAT'' trace is not
an $R$-derivation of $\bot$ until \emph{soundness of the trace} has been certified inside $R$. Consequently a
superpolynomial EF lower bound would give only ``P$\ne$NP is consistent with PV''. That is OPEN, and the
recent literature conjectures it but has not proved it.

\section{Computational data (COMPUTATIONAL)}
The script \texttt{worst-case-derivation/treelike\_min.py} computes exact minimum tree-like resolution size,
minimizing over all refutations. It uses the exact correspondence between tree-like refutations and decision
trees: refutation size is $2L-1$ for a tree with $L$ leaves. Each instance is a single run under 1\,s, so no
chunking was needed. Results are in \texttt{results.jsonl}:
\begin{center}\begin{tabular}{lrrr}
instance & vars & clauses & min tree-like size\\\hline
$\mathrm{PHP}^2_1$&2&3&5\\ $\mathrm{PHP}^3_2$&6&9&21\\ $\mathrm{PHP}^4_3$&12&22&85\\
Tseitin$(K_4,\text{odd})$&6&16&55
\end{tabular}\end{center}
These sizes are consistent with the known $2^{\Theta(n\log n)}$ tree-like bound for PHP. They illustrate
the consequent side only, which Theorem~\ref{t:reform} shows carries no information about the bridge.

\section{Role assignments (CANDIDATE work items)}
\begin{enumerate}
\item \textbf{Lower-bound lead:} replace the bridge with the PV-relative form (Prop.~\ref{p:pv}). Target:
for the extension-enabled hypergraph calculus $R_{\mathrm{hyp}}$, show $R_{\mathrm{hyp}}$ p-simulates EF and
vice versa, so that Cook's theorem transfers exactly.
\item \textbf{Reduction builder:} prove or refute that $R_{\mathrm{hyp}}$ without extension is p-equivalent to
resolution. If it is, Theorem~\ref{t:reform} applies and the route is a reformulation.
\item \textbf{Hard instances:} only candidates for \emph{EF-hard} families matter, for example
$\neg$Circuit-lower-bound or random $\Delta$-CNF families. Note the Pich--Santhanam (FOCS 2019) barrier:
under standard assumptions, lower bounds for natural distributions of hard tautologies are infeasible to
prove.
\item \textbf{Resource auditor:} certify that $\operatorname{cost}_R$ charges extension definitions and
hyperedge size polynomially. A mischarge collapses $\overline W$.
\item \textbf{Red team:} for every proposed bridge proof, check whether it relativizes. By
Theorem~\ref{t:rel}, any relativizing argument is automatically wrong.
\end{enumerate}

\section*{CORPUS REVIEWED / DUPLICATION AUDIT}
Reviewed: \texttt{AGENTS.md}, \texttt{research/agents/claude/*} (world ledger L13, L34), and a repository
search for $\overline W$, ``derivation cost'' and ``weakest sufficient''. Hits are in LSPV-1/5 reports, for
example \texttt{lspv-1/2026-11-03\_1800\_..._376.tex} and \texttt{lspv-1/..._LOGDEGREE\_RESPC\_TSEITIN479.md},
which treat resolution and RES-PC Tseitin costs. Those reports concern the consequent for weak systems.
Theorem~\ref{t:reform} explains why they cannot close the bridge. No contradiction was found. Root
\texttt{main.tex} was not edited, because this routine owns no ownership region in it. \textbf{Stale or
overstated claims:} none were identified in the files read.

\section*{References}
Cook (1975) STOC; Cook--Reckhow (1979) JSL; Haken (1985) TCS; Urquhart (1987) JACM;
Kraj\'{\i}\v{c}ek, \emph{Bounded Arithmetic, Propositional Logic and Complexity Theory} (1995);
Pich--Santhanam, ``Why are proof complexity lower bounds hard?'' FOCS 2019;
Oliveira et al., ``Meta-mathematics of computational complexity theory'' arXiv:2506.16956 (2025);
Atserias--M\"uller, ``Automating resolution is NP-hard'' (2019).
\end{document}
% END_OF_VERIFIED_REPORT::CLAUDE::WCDERIV-0001
